\begin{abstract}
Part 1 asks for five stopping-reason labels on 1,171 terminated trials, only 157 of them labeled.
K-means diagnostics show weak cluster structure in the reason embeddings (silhouette 0.16--0.31, adjusted Rand index with the true label-sets at most 0.20).
We compared nine labeling approaches, from cluster voting to a semi-supervised GMM, self-training and label spreading.
The selected model averages a supervised logistic regression (LR) with two pseudo-label models and raises validation sample-F1 from 0.865 to 0.897 and macro-F1 from 0.707 to 0.806.
Most of the macro-F1 gain comes from Funding, which has 8 labeled positives.
Part 2 predicts whether a trial met its primary endpoint in each phase.
The starter model's log loss of 0.940 is worse than predicting each phase's success rate, because its probabilities are poorly calibrated.
One regularized LR on field-wise TF-IDF features, shared across phases by feature augmentation, lowers mean validation log loss to 0.567 and raises ROC-AUC from 0.690 to 0.761.
A split by start year suggests that temporal drift costs about 0.07 AUC and leaves log loss almost unchanged.
\end{abstract}
